\chapter{Các hàm Python bổ trợ và giải hệ tuyến tính}
\label{app:api}

\section{Đọc đúng tên hàm và thứ tự hệ số}
Các hàm dưới đây nằm trong \texttt{pps/chapter1\_tools.py}; có thể gọi bằng
\begin{verbatim}
from pps import chapter1_tools as ch1
\end{verbatim}
Chúng không phải tất cả đều có nút riêng trong GUI. Đặc biệt \textbf{hệ số đơn thức của mô-đun này theo bậc tăng}; ô Horner trên GUI lại theo bậc giảm. Ví dụ $2x^3-3x+1$ là \texttt{[1,-3,0,2]} khi gọi \texttt{horner\_division}, và \texttt{2 0 -3 1} khi nhập GUI.

\begin{center}\small\begin{tabular}{@{}p{.38\textwidth}p{.54\textwidth}@{}}\toprule
API & Vai trò và nơi đọc thuật toán\\\midrule
\texttt{chebyshev\_value(n,z)} & Tính $T_n(z)$ bằng truy hồi; xem ngay dưới bảng\\
\texttt{horner\_division(c,a)} & Chia đa thức đơn thức cho $x-a$; Chương \ref{ch:interp}\\
\texttt{newton\_horner(nodes,c,z)} & Horner dạng Newton theo đúng thứ tự mốc; Chương \ref{ch:interp}\\
\texttt{newton\_append(nodes,c,x,y)} & Thêm một mốc mà không tính lại hệ số cũ; Chương \ref{ch:interp}\\
\texttt{interpolation\_error\_bound} & Cận tại một điểm khi đề cho $M_m$; Chương \ref{ch:interp}\\
\texttt{interpolation\_global\_bound} & Cận toàn dải lưới đều; cần tên mẫu \texttt{'cách đều'}\\
\texttt{chebyshev\_global\_bound} & Cận trên $[a,b]$ cho đúng $m$ nghiệm Chebyshev; không lấy min/max mốc thay $a,b$\\
\texttt{fill\_missing} & Điền ô thiếu theo giả thiết bậc đa thức; quy trình dưới đây\\
\bottomrule\end{tabular}\end{center}

\begin{center}\small\begin{tabular}{@{}p{.38\textwidth}p{.54\textwidth}@{}}\toprule
API (tiếp) & Vai trò và nơi đọc thuật toán\\\midrule
\texttt{scan\_level}, \texttt{monotone\_branches} & Quét mức và nhánh đơn điệu của \emph{bảng}; Chương \ref{ch:inverse}\\
\texttt{inverse\_by\_swap} & Wrapper hàm ngược; dùng đúng tập mốc đã chọn\\
\texttt{bisect\_model} & Chia đôi một mô hình đã có, trả lịch sử; Chương \ref{ch:inverse}\\
\texttt{inverse\_controlled} & Wrapper lặp có chứng nhận, mặc định tối đa 1000 vòng; GUI mặc định 100\\
\texttt{spline\_degree} & API cũ bậc 1--4, biết $p-1$ đạo hàm trái; Chương \ref{ch:spline}\\
\texttt{spline\_boundary} & API cũ bậc 1--4, đúng $p-1$ điều kiện biên; GUI mới xét cả hệ thiếu/thừa\\
\texttt{spline\_at} & Giá trị/đạo hàm, kiểm tra hai phía tại nút; không tự ngoại suy\\
\texttt{thomas} & Giải hệ ba đường chéo có đổi thang từng hàng; ví dụ dưới đây\\
\texttt{gaussian\_elimination} & Hệ vuông, đổi thang hàng và chọn trụ theo cột; ví dụ dưới đây\\
\texttt{transformed\_fit} & Bốn mô hình biến đổi cố định; Chương \ref{ch:ls}\\
\texttt{original\_scale\_fit} & API cũ gọi tối ưu SciPy; khác nhánh GUI LM có đầy đủ vết bước thử\\\bottomrule
\end{tabular}\end{center}

Một \texttt{Report} thường chứa \texttt{value}, \texttt{lines}, \texttt{tables}. Tuy nhiên kiểu \texttt{value} phụ thuộc hàm: vector nghiệm, số, đối tượng spline hoặc hàm dự đoán. Không giả sử mọi kết quả đều có cùng cấu trúc.

\section{Truy hồi Chebyshev và cận sai số}
\[
 T_0(z)=1,\quad T_1(z)=z,\quad
 T_{k+1}(z)=2zT_k(z)-T_{k-1}(z).
\]
Tính theo $k=1,\ldots,n-1$, giữ hai giá trị cũ trước khi gán cặp mới. Ví dụ $z=0.5$: $T_0=1,T_1=.5,T_2=-.5,T_3=-1$. Hàm chấp nhận $z$ ngoài $[-1,1]$ nhưng có thể tràn vì độ lớn tăng nhanh. $n$ phải nguyên không âm.

Với $m$ mốc và cận đạo hàm cấp $m$ là $M_m$:
\[
 |f(z)-P(z)|\le\frac{M_m}{m!}\prod_{i=0}^{m-1}|z-x_i|.
\]
$M_m$ phải đúng trên đoạn chứa mọi mốc và điểm tính, kể cả khi ngoại suy. Với $m\ge2$ mốc đều bước $h$, cận toàn dải mà API dùng là
\[
 |f-P|\le\frac{M_m h^m}{4m}.
\]
Với $m$ nghiệm Chebyshev trên $[a,b]$, cận là
\[
 \frac{2M_m}{m!}\left(\frac{b-a}{4}\right)^m.
\]
Ví dụ $x=(0,1,2)$, $z=.5$, $M_3=6$: cận điểm bằng $.375$; cận lưới đều bằng $.5$. Với 3 nghiệm Chebyshev trên $[0,2]$ cùng $M_3$, cận bằng $.25$. Đây là cận khi có giả thiết đạo hàm, không phải sai số thực tế suy từ bảng đo.

\section{Điền ô thiếu khi đề cho trước bậc}
\texttt{fill\_missing(x,y,degree,tolerance=1e-9)} nhận các ô thiếu dạng \texttt{None} hoặc \texttt{nan}. Quy trình:
\begin{enumerate}
\item Kiểm tra hoành độ hữu hạn, phân biệt, đủ $d+1$ tung độ biết; $0\le d<m$.
\item Chọn $d+1$ ô biết đầu tiên theo thứ tự mảng để dựng đa thức Newton bậc không quá $d$.
\item Tính giá trị mô hình ở \emph{mọi} hoành độ, kể cả các ô biết còn lại.
\item Với các ô biết, kiểm $\max|P(x_i)-y_i|\le\tau\max|y_i|$. Nếu thang tung độ bằng 0 thì vế phải là 0; không tự thêm một ngưỡng tuyệt đối 1.
\item Chỉ điền các ô thiếu; giữ nguyên các ô đã biết nếu phép kiểm đạt.
\end{enumerate}
Ví dụ $x=(0,1,2,3)$, $y=(1,\mathrm{None},5,10)$, $d=2$. Dùng ba cặp biết $(0,1),(2,5),(3,10)$: tỷ sai phân cấp 1 là $2,5$, cấp 2 là $(5-2)/(3-0)=1$. Vì vậy
\[
 P(x)=1+2x+x(x-2)=1+x^2,\qquad y(1)=2.
\]
Kết quả $(1,2,5,10)$. Nếu có thêm một ô biết không nằm trên đa thức trong ngưỡng kiểm thì hàm báo mâu thuẫn, không sửa ô đó để ép giả thiết.

\section{Thomas: thứ tự khử tiến và thế ngược}
Nhập bốn vector cùng độ dài $n$: đường chéo dưới $a$, chính $b$, trên $c$ và vế phải $d$, với $a_0=c_{n-1}=0$. Hệ là
\[
 a_ix_{i-1}+b_ix_i+c_ix_{i+1}=d_i.
\]
API chia mỗi hàng cho $s_i=\max(|a_i|,|b_i|,|c_i|)$ trước khi khử; hàng toàn 0 bị từ chối. Các công thức dưới đây áp dụng cho hệ \emph{sau đổi thang}:
\begin{align*}
 m_i&=a_i/\widetilde b_{i-1},\\
 \widetilde b_i&=b_i-m_i c_{i-1},\qquad
 \widetilde d_i=d_i-m_i\widetilde d_{i-1},\quad i=1,\ldots,n-1,\\
 x_{n-1}&=\widetilde d_{n-1}/\widetilde b_{n-1},\\
 x_i&=(\widetilde d_i-c_ix_{i+1})/\widetilde b_i,\quad i=n-2,\ldots,0.
\end{align*}
Thomas không chọn đổi hàng. Pivot gần 0 có thể xuất hiện ở một hệ vẫn khả nghịch; khi đó cần bộ giải có chọn trụ. Dung sai pivot mặc định $10^{-12}$ xét trên hệ đã đổi thang. Không gọi lỗi pivot là một bằng chứng hệ gốc vô nghiệm.

\subsection{Ví dụ trọn vẹn ba ẩn}
\[
 \begin{bmatrix}2&-1&0\\-1&2&-1\\0&-1&2\end{bmatrix}x
 =\begin{bmatrix}1\\0\\1\end{bmatrix}.
\]
Chia cả ba hàng cho 2: $b=(1,1,1)$, $a=(0,-.5,-.5)$, $c=(-.5,-.5,0)$, $d=(.5,0,.5)$.
\begin{center}\begin{tabular}{rrrrr}\toprule
$i$&$m_i$&$\widetilde b_i$&$\widetilde d_i$&Nguồn vế phải\\\midrule
0&---&1&$1/2$&Khởi tạo\\
1&$-1/2$&$3/4$&$1/4$&$0-(-1/2)(1/2)$\\
2&$-2/3$&$2/3$&$2/3$&$1/2-(-2/3)(1/4)$\\\bottomrule
\end{tabular}\end{center}
Thế ngược lần lượt $x_2=1$, $x_1=(1/4+1/2)/(3/4)=1$, $x_0=(1/2+1/2)/1=1$. Kiểm hệ gốc cho phần dư $(0,0,0)$.
\begin{verbatim}
r = ch1.thomas([0,-1,-1], [2,2,2], [-1,-1,0], [1,0,1])
print(r.value)  # [1. 1. 1.]
\end{verbatim}

\section{API giải hệ: input và output}
\texttt{gaussian\_elimination(A,b,pivot\_tolerance=1e-12)} cần $A$ vuông hữu hạn, $b$ cùng cỡ. Không dùng API này để giải trực tiếp bài chữ nhật $\Phi c\approx z$; dùng nhánh hồi quy và các cách giải ở Chương \ref{ch:ls}.
Output là vector nghiệm, bảng chọn pivot thực tế và phần dư theo hệ gốc. Phần này không nhắc lại thuật toán khử.
Hạng thiếu trong dung sai bị từ chối; hàm không phân loại riêng vô nghiệm/vô số nghiệm như bộ báo cáo spline tổng quát. Nhật ký API này lưu bảng chọn pivot, không lưu mọi ma trận sau từng phép khử.

\section{Phân biệt API cũ và luồng lời giải hiện tại}
\texttt{spline\_degree} trả \texttt{PPoly}, truyền đạo hàm đầu trái cho bậc 1--4. \texttt{spline\_boundary} cũ nhận đúng $p-1$ điều kiện; GUI mới qua \texttt{solve\_spline} kiểm cả hạng và điều kiện dư. Không dùng giới hạn của API cũ để kết luận GUI chỉ làm được biên trái.

\texttt{original\_scale\_fit} cũ gọi \texttt{scipy.optimize.least\_squares}, với tối đa 10\,000 lần đánh giá và các ngưỡng nội bộ $10^{-12}$. Nó không cung cấp nhật ký LM của GUI. Nhánh GUI \textbf{SSE thang y} dùng \texttt{solve\_least\_squares}, ghi bước thử, tham số, hệ Jacobian, hệ số giảm chấn và bước nhận/từ chối. Hai cách không nhất thiết có cùng số vòng, cùng quỹ đạo hay cùng tiêu chuẩn dừng.

Các wrapper nội suy ngược cũ gọi bộ giải mới nhưng chuyển trạng thái thất bại thành ngoại lệ qua \texttt{require\_success}. Đọc báo cáo của \texttt{solve\_inverse} nếu cần giữ lịch sử thất bại để kiểm tra. Không lấy giá trị cuối của một phép lặp thất bại làm đáp số được công nhận.
